\documentclass{situationist-preprint}

\title{Queerness Is Not Evidence}
\msid{MTE-05-b}
\msfamily{MTE-05}
\msdiscipline{Metaethics}
\mstype{Companion manuscript}
\mskeywords{error theory; argument from queerness; Mackie; evidence; companions in guilt}

\begin{document}
\maketitle

\begin{abstract}
MTE-05-a argues that Mackie's argument from queerness proves too much, since mathematical facts bear all three marks of queerness. A defender of the error theory may reply that queerness is not offered as a proof but as evidence: it lowers the probability of moral facts without excluding them. We argue that queerness is not evidence either. Queerness is a property of a hypothesis relative to a background theory, not an observation that the hypothesis predicts less well than its rival. What remains of the argument is a consideration of parsimony, and parsimony counts against moral facts only if they do no explanatory work, which is a different argument.
\end{abstract}

\section{The evidential reading}\label{s:ev}

\citet[ch.~1.9]{Mackie1977} presents the argument from queerness as an argument against the existence of objectively prescriptive facts. MTE-05-a takes its principle to be that whatever is queer does not exist (Principle P1 of MTE-05-a) and shows that the principle would make all positive mathematical claims false (Proposition~4.1 and Theorem~4.2 of MTE-05-a). Some error theorists read the argument more modestly, as a consideration that weighs against moral facts and could be outweighed \citep{Olson2014}. On this reading, the companions in guilt reply shows only that mathematics has a counterweight that morality lacks.

The evidential reading requires that queerness be evidence. An observation $E$ is evidence against a hypothesis $H$, relative to a rival $H'$, when $E$ is less probable on $H$ than on $H'$. The queerness of moral facts is not such an observation. It is not something we observe and that the hypothesis of moral facts predicts poorly; it is a description of the hypothesis itself, namely that the facts it posits are unlike the others in our inventory. The hypothesis that there are moral facts and the hypothesis that there are none do not differ in how probable they make the fact that moral facts, if there are any, are unlike tables and electrons. Both entail it, if the marks are correctly stated. So the likelihood ratio is one, and queerness is not evidence.

\begin{objection}
Strangeness relative to a successful background theory is a reason for a low prior, and a low prior is all the error theorist needs.
\end{objection}

\begin{reply}
A prior lowered by unlikeness to our other commitments is a prior lowered by parsimony: posit no more kinds than one needs. Parsimony counts against a kind only if it does no explanatory or justificatory work. Whether moral facts do such work is the question of whether they figure in our best explanations of actions, institutions and moral beliefs \citep{Sturgeon1985}, and whether our reasons for belief, which are themselves normative, can be had without them \citep{Cuneo2007}. That is an argument of a different kind, whose force does not come from queerness. The history of science, moreover, shows that unlikeness to the current inventory is a poor guide to existence: fields acting at a distance and non-local correlations were queer by the standards of the theories they displaced.
\end{reply}

\section{Conclusion}

Read as proof, the argument from queerness proves too much, as MTE-05-a shows. Read as evidence, it proves nothing, since queerness describes the hypothesis rather than the data. Since moral facts, by MTE-01, are natural facts about idealized reactions, the marks do not even apply to them as Mackie states them. Not all moral talk is false.

\begin{thebibliography}{99}\small
\bibitem[Cuneo(2007)]{Cuneo2007} Cuneo, T. (2007). \emph{The Normative Web: An Argument for Moral Realism}. Oxford University Press.
\bibitem[Mackie(1977)]{Mackie1977} Mackie, J.~L. (1977). \emph{Ethics: Inventing Right and Wrong}. Penguin.
\bibitem[Olson(2014)]{Olson2014} Olson, J. (2014). \emph{Moral Error Theory: History, Critique, Defence}. Oxford University Press.
\bibitem[Sturgeon(1985)]{Sturgeon1985} Sturgeon, N.~L. (1985). Moral explanations. In D. Copp and D. Zimmerman (eds.), \emph{Morality, Reason and Truth}. Rowman and Allanheld.
\end{thebibliography}

\end{document}
